\chapter*{TÓM TẮT}
\addcontentsline{toc}{chapter}{TÓM TẮT}
\fontsize{13}{15}\selectfont

\begin{sloppypar}
Đồ án với đề tài \textit{"Thiết kế hệ thống AIoT phát hiện và cảnh báo bất thường cho máy giặt dân dụng sử dụng Autoencoder không giám sát trên thiết bị biên"} tập trung xây dựng hệ thống AIoT phát hiện và cảnh báo bất thường cho máy giặt dân dụng. Hệ thống tích hợp cảm biến âm thanh INMP441 và cảm biến gia tốc ADXL345, xử lý dữ liệu trực tiếp trên vi điều khiển Seeed Studio XIAO ESP32-S3 theo kiến trúc ba trạng thái cho ba pha vận hành GENTLE, STRONG và SPIN.
\end{sloppypar}

Về mặt thuật toán, đồ án xây dựng mô hình tự mã hóa đối xứng, kết hợp trích xuất đặc trưng tín hiệu số, lượng tử hóa INT8 và cơ chế phân tách mô hình theo từng pha vận hành. Cách tiếp cận này giúp hệ thống học quy luật của trạng thái bình thường, sau đó phát hiện bất thường thông qua sai số tái tạo mà không cần tập dữ liệu lỗi được gán nhãn đầy đủ.

Về mặt triển khai, hệ thống được hoàn thiện theo mô hình từ biên đến ứng dụng. Kết quả suy luận tại vi điều khiển Seeed Studio XIAO ESP32-S3 được truyền qua MQTT trên HiveMQ Cloud, lưu lịch sử bằng Firebase Firestore và hiển thị trên ứng dụng di động Flutter; nhờ đó người dùng có thể theo dõi trạng thái máy giặt và nhận cảnh báo gần thời gian thực.

Kết quả thực nghiệm cho thấy mô hình đạt F1-Score lần lượt 0,9898, 0,9892 và 0,9967 cho ba pha GENTLE, STRONG và SPIN; AUC đều lớn hơn 0,98. Độ trễ suy luận trung bình trên vi điều khiển Seeed Studio XIAO ESP32-S3 đạt 6,4 ms, dữ liệu đo xa được gửi với tần suất 1 Hz, độ trễ từ biên đến ứng dụng dưới 2 giây, trong khi firmware hoàn chỉnh sử dụng khoảng 1,20 MB Flash và 62 KB SRAM. Các kết quả này cho thấy giải pháp có khả năng vận hành ổn định trên phần cứng tài nguyên thấp và phù hợp để tích hợp vào các hệ thống giám sát máy giặt dân dụng.

\vspace{0.5cm}
\textbf{\textit{Từ khóa:}} \textit{AIoT, TinyML, Mạng nơ-ron tự mã hóa, Phát hiện bất thường, ESP32-S3.}
